% !TeX root = ../thesis.tex
\chapter{Introduction}
\label{ch:introduction}

\section{Positron Emission Tomography: Principles and Practice}
\topics{annihilation coincidence detection in one paragraph; PET as functional/molecular
imaging versus the anatomical modalities (CT, MRI); clinical roles (oncology staging,
neurology, cardiology) and preclinical research roles; hybrid PET/CT and PET/MRI as the
modern norm; the basic scanner anatomy: scintillator, photodetector, electronics,
reconstruction.}

Positron Emission Tomography (PET) is a nuclear imaging modality that enables the non-invasive
visualisation of biological processes in vivo. It relies on the detection of pairs of gamma 
photons produced by the annihilation of positrons emitted from radiotracers administered
into the subject. PET radiotracers are generally bound to specific biological targets, This differs from anatomical imaging modalities like CT and MRI, 
which provide structural information. PET is widely used in clinical settings for 
oncology, neurology, and cardiology, as well as in preclinical research to study
disease mechanisms and evaluate new therapies.

\section{Cost Drivers in PET Systems}
\topics{scintillator crystal cost per cm\textsuperscript{3} (LYSO vs BGO) and its share of
the bill of materials; photodetector and readout channel count scaling; precision
fabrication of pixellated arrays (dicing, polishing, reflector assembly); radiotracer
supply chain and cyclotron proximity; siting, shielding and staffing; service and total
cost of ownership; how these costs translate into scan price and limited scanner
availability.}

\section{Opportunities for Low-Cost PET}
\topics{cheaper scintillator materials (BGO revival, plastics, nanocomposites, ceramics);
monolithic slabs eliminating dicing and dead space; reduced or sparse detector coverage
and partial-ring/flat-panel geometries; multiplexed or low-channel-count readout;
software and ML compensating for cheaper hardware; sensitivity--cost trade-offs
highlighted by total-body PET; open hardware and academic scanner efforts.}

\section{Aims of This Review}
\topics{map the whole PET pipeline with a cost lens; identify which cost--performance
trade-off levers are under-explored; establish evaluation criteria for candidate
directions (cost impact, novelty, publishability, feasibility of an experimental
demonstrator within a PhD); converge on a defensible project direction.}

\section{Candidate Research Directions}
\topics{continuous-coordinate / geometry-agnostic reconstruction (prior pilot work);
flat-panel and box geometries from monolithic slabs; nanocomposite or ceramic
scintillator characterisation; neural event localisation and the simulation-to-real
transfer problem; low-channel-count readout schemes; how each scores against the
supervisor's constraints: cheaper PET, personal interest, publishable, extensible to a
real-world experimental model.}

\section{Thesis Outline}
\topics{one sentence per chapter; note that the document currently comprises the
introduction and literature review, with experimental chapters to follow once the
direction is fixed.}
